% problem_id: S046-thm:lequiv-contr-hae
% source_id: S046 (HoTT Book)
% source_file: equivalences.tex
% statement_locator: equivalences.tex lines 447--449
% proof_locator: equivalences.tex lines 450--458
% env: thm
% pairing: tight (verified)
% status: complete (statement + author proof verbatim + reason + locators)
%
\begin{thm}\label{thm:lequiv-contr-hae}
For any $f:A\to B$ we have ${\iscontr(f)} \to {\ishae(f)}$.
\end{thm}

\begin{proof}
Let $P : \iscontr(f)$. We define an inverse mapping $g : B \to A$ by sending each $y : B$ to the center of contraction of the fiber at $y$:
\[ g(y) \defeq \proj{1}(\proj{1}(Py)). \]
We can thus define the homotopy $\epsilon$ by mapping $y$ to the witness that $g(y)$ indeed belongs to the fiber at $y$:
\[ \epsilon(y) \defeq \proj{2}(\proj{1}(P y)). \]
It remains to define $\eta$ and $\tau$. This of course amounts to giving an element of $\rcoh{f}{g}{\epsilon}$. By \cref{lem:coh-hfib}, this is the same as giving for each $x:A$ a path from $(gfx,\epsilon(fx))$ to $(x,\refl{fx})$ in the fiber of $f$ over $fx$. But this is easy: for any $x : A$, the type $\hfib{f}{fx}$
is contractible by assumption, hence such a path must exist. We can construct it explicitly as
\[\opp{\big(\proj{2}(P(fx))(gfx,\epsilon(fx))\big)} \ct \big(\proj{2}(P(fx)) (x,\refl{fx})\big). \qedhere \]
\end{proof}
